\documentclass{ximera}
\title{Directional Derivatives CW}

\newcommand{\pskip}{\vskip 0.1 in}

\begin{document}
\begin{abstract}
Introduction to directional derivatives.
\end{abstract}
\maketitle

\section{Introduction}

We've learned how to compute rate of change of a temperature function $T=f(x,y)$ in the directions of the vectors ${\bf i} = \langle 1, 0\rangle$, and ${\bf j} = \langle 0, 1 \rangle$. This section is about computing the rate of change of the temperature function (with respect to distance) in the direction given by a vector ${\bf v} = \langle v_1 ,v_2 \rangle$. 

The idea is to approximate the temperature function with a linear function in the neighborhood of a point. So we'll start by assuming the temperature function is linear.

\emph{Note:} To compute directional derivatives, we need not use a temperature function, but it is necessary that the inputs of the function have the same units..


\section{Linear Functions}

\begin{question}  \label{Ee53efehhy}

Some level curves of a temperature function $T = f(x,y)$ are shown below. The output is measured in Celsius degrees and the input in meters.

\pdfOnly{
Access Geogebra interactives through the online version of this text at
 
\href{https://www.geogebra.org/classic/htee9xke}.
}
 
\begin{onlineOnly}
    \begin{center}
\geogebra{htee9xke}{900}{600}
\end{center}
\end{onlineOnly}

Access Geogebra interactives through the online version of this text at
 
\href{https://www.geogebra.org/classic/htee9xke}{163: Directional Derivative}.


\begin{enumerate}
\item Find an expression for $f$.
\[
     T = f(x,y) = \answer{2x + 4y}
\]

\item Use the sliders $\theta$ and $s$ to approximate the (instantaneous) rate at which the temperature changes from the origin in the directions of the following vectors. Include units.

Note:

\begin{itemize}

\item The slider $\Delta s$ controls the distance (in meters) between consecutive tick marks on the number line through the origin. 

\item The slider $k$ controls the temperature of the red level curve. That is, the red level curve has equation $f(x,y)=k$.

\item The slider $\Delta k$ controls the spacing of the level curves. That is, $\Delta k$ is the temperature difference between consecutive level curves.  

\end{itemize}

\begin{enumerate}

\item ${\bf v} = \langle 1 ,2 \rangle$, where $\theta \sim 1.1$

\item ${\bf w} = \langle -3,1\rangle$, where $\theta \sim 2.8$.

\end{enumerate}

\item Use a  difference quotient to compute the exact values of the above rates of change. Then compare these with your estimates. Click the arrow to the lower right for a partial solution.

\begin{expandable}
Because the temperature function is linear, we can avoid limits and compute the rate of change at the origin in the direction of the vector ${\bf v} = \langle 1, 2\rangle$ as a difference quotient.

In moving from the origin to the point $(1,2)$ the average rate at which the temperature changes with respect to the distance travelled is
\[
   \frac{\Delta T}{\Delta s} = \frac{f(1,2) - f(0,0)}{|\langle 1, 2 \rangle|}\frac{^\circ C}{m} = \frac{10}{\sqrt{5}}\frac{^\circ C}{m} .
\]

Then make a similar computation for rate of change in the direction of the vector ${\bf w} = \langle -3,1\rangle$.
\end{expandable}


\item Use a difference quotient to compute the rate  at which the temperature changes from the origin in the direction of the vector ${\bf v}=\langle v_x , v_y \rangle$. 

\item Check that your expression in the previous part has the correct units.

\item Interpret your expression for the rate of change as the scalar projection of one vector onto another. Then write this scalar projection as the product of a magnitude of the gradient vector and a trigonometric function of the angle between the gradient vector and the vector ${\bf v}$.

\end{enumerate}
\end{question}




\end{document}